\section{Methodology}
\label{sec:method}

\para{Models.}
We study three pretrained Pre-LN decoder-only Transformers loaded in FP32 on
a single NVIDIA RTX A6000 (49 GB): \gptsmall (124M parameters), \gptmedium
(355M), and \pythia (410M). \gptsmall is the canonical baseline used by
\citet{brody2023expressivity}; \gptmedium adds a scale check; \pythia
introduces a distinct training corpus and is the model used by
\citet{gupta2024geometric}. All checkpoints are pulled from HuggingFace at
fixed revisions; random seeds are set globally across \texttt{random},
\texttt{numpy}, and \texttt{torch}.

\subsection{LayerNorm Intervention Recipes}
\label{sec:method:recipes}

We hook every \verb|nn.LayerNorm| in each Transformer block and replace its
forward pass with one of the recipes in \tabref{tab:recipes}. Each non-stock
recipe is applied to one of four layer ranges: \tgtall (every block),
\tgtearly (blocks $0\dots\lfloor n/3\rfloor$), \tgtmid
($\lfloor n/3\rfloor\dots\lfloor 2n/3\rfloor$), and \tgtlate
($\lfloor 2n/3\rfloor\dots n$). The final-block \layernorm and the LM-head
\layernorm are always preserved, following \citet{kanavalau2026taper}'s
finding that the final normalisation is the critical scale anchor.

\begin{table}[t]
\centering
\caption{\layernorm intervention recipes. $\mu$ and $\sigma$ are the
per-token mean and standard deviation; $\gamma,\beta$ are the learned
affine parameters.}
\label{tab:recipes}
\begin{tabular}{@{}lll@{}}
\toprule
\textbf{Recipe} & \textbf{Replacement} & \textbf{Geometric meaning} \\
\midrule
\stockln    & $\layernorm(\vx)$ (control) & --- \\
\nomean     & $\vx \cdot \text{rsqrt}(\overline{\vx^2}) \cdot \gamma + \beta$ & \rmsnorm-equivalent \\
\noaffine   & $(\vx-\mu)/\sigma$ & drop $\gamma,\beta$ \\
\weakened   & $\alpha\cdot\layernorm(\vx)+(1-\alpha)\vx$, $\alpha\in\{0.25,0.5,0.75\}$ & partial \layernorm \\
\identityln & $\vx$ & bypass entirely \\
\bottomrule
\end{tabular}
\end{table}

\subsection{Divergent Association Task}
\label{sec:method:dat}

We adopt the \dat protocol of \citet{olson2021naming} and
\citet{chen2023probing}: produce $10$ unrelated nouns, retain $7$ that are
valid \glove-300d entries, and score
\[
\dat = \frac{100}{n(n-1)} \sum_{i\neq j} \big(1 - \cos(\vv_i, \vv_j)\big),
\]
with \glove-6B-300d \citep{pennington2014glove}.

\para{Iterative one-noun protocol.}
Base \gptsmall cannot reliably emit a $10$-noun list in a single shot. We
therefore use an iterative protocol: anchored by a random seed word from
$\{\text{apple, shoe, river, engine, book, wallet, candle, fence, mirror,
rocket}\}$, the model is asked to emit one new noun at a time conditioned on
the growing list. This produces clean, parseable single-token nouns at every
step. Absolute \dat numbers should be interpreted under this protocol;
relative comparisons across interventions are within-protocol and valid.

\para{Validity gating.}
For each $10$-slot generation we record a validity rate: the fraction of
slots that produced a \glove-valid English noun. We require validity
$\ge 0.9$ to call a measurement ``coherent''. Below this threshold,
generations are gibberish and \dat becomes uninterpretable: out-of-vocabulary
tokens have effectively random \glove embeddings and therefore score
spuriously high (we document this artifact in \secref{sec:results:96trio}).

\subsection{Decoding Strategies}
\label{sec:method:decoding}

We sweep three decoding regimes:
(i) \greedy: deterministic argmax;
(ii) \topp: top-$p{=}0.9$, $T{=}0.7$ \citep{holtzman2020curious};
(iii) \hight: $T{=}1.2$, top-$p{=}0.95$.
All three set \verb|repetition_penalty|$=1.6$ and
\verb|no_repeat_ngram_size|$=2$ to suppress \gptsmall's degenerate
single-token loops.

\subsection{Per-layer Concept-Association Geometry}
\label{sec:method:geometry}

For a fixed bank of $20$ unrelated nouns, we capture last-token hidden states
at every layer and compute three quantities: (i) mean pairwise cosine
distance over all $\binom{20}{2}=190$ pairs (concept distinguishability),
(ii) the L2-norm distribution std/mean ratio (norm-collapse evidence
following \citet{gupta2024geometric}), and (iii) the angle to the uniform
vector $\bm{1}$ (sanity check on \layernorm's projection).

\subsection{Statistics}
\label{sec:method:stats}

Because \dat distributions are skewed and sample sizes are modest, we use
non-parametric tests. The primary test is a Mann--Whitney $U$ test of \dat
under each intervention against \stockln. We report Cohen's $d$ as the
effect size. For per-layer geometry we use a Wilcoxon paired test with each
of the $190$ pairs as the matched unit. Multiple-comparison correction uses
Bonferroni at $\alpha{=}0.05/18{\approx}0.003$ over the $18$
fluency-preserving \nomean cells.

\subsection{Sample sizes and compute}
\label{sec:method:samples}

Each (model, decoding, intervention) cell is run with $20$ generations,
each producing $10$ noun-slots --- $13$ cells $\times\,20\,\text{runs}\times
10\,\text{slots}=2{,}600$ noun generations per (model, decoding) pair. The
high-$T$ panel uses $n{=}15$ runs per cell. Wall time is approximately
$7\!-\!15$ minutes per (model, decoding) for the full \dat grid;
$\sim\!30\,\text{s}$ for geometry and $\sim\!10\,\text{s}$ for perplexity.
All raw outputs (per-step generated nouns), per-layer hidden-state
statistics, and configuration JSONs are persisted to \texttt{results/}.
